\section{Discussion}
\label{sec:discussion}

Our empirical findings and architectural analysis yield several broader insights for the design of long-running autonomous AI agents.

\subsection{Why Versioning Matters for AI Agents}
The dominant paradigm in commercial and open-source agent memory frameworks has treated memory as a mutable database table~\cite{packer2023memgpt,deshmukh2024mem0}. When an agent updates a belief, it issues an in-place overwrite. Our results demonstrate that this design is fundamentally flawed for three reasons:
\begin{enumerate}
    \item \textbf{Loss of Auditing Lineage:} In regulated domains (e.g., healthcare, automated financial trading, legal analysis), an agent's reasoning must be forensically reproducible. Destructive updates destroy the historical audit trail. With ManiGit's immutable commit DAG, an auditor can reconstruct the agent's exact operative context at any past moment in time.
    \item \textbf{Inability to Answer Retrospective Queries:} Users routinely ask agents questions that hinge on temporal context (e.g., ``What was our configuration before yesterday's migration?''). Systems without versioning fail completely on these queries, as verified by our 0\% baseline historical recovery results.
    \item \textbf{Branch Isolation for Speculative Reasoning:} Autonomous agents frequently need to explore alternative hypotheses or test conflicting execution plans. Git-like branching allows agents to fork their context into isolated branches, test outcomes, and abandon speculative branches without contaminating the primary operational state.
\end{enumerate}

\subsection{Decoupling Memory from Context}
A core conceptual contribution of this work is establishing that \textbf{Memory $\ne$ Context}. 

In early LLM research, memory was synonymous with the prompt buffer. Expanding context windows to hundreds of thousands of tokens led many practitioners to assume that external memory architectures were obsolete: one could simply append the entire interaction history to the prompt.

Our experiments refute this assumption. As demonstrated by Full History's performance at $B=128$ and $B=256$, relying on raw conversational history is vulnerable to context budget exhaustion. When histories grow long, older critical constraints are severed. Even when budgets are large ($B=2048$), Full History consumes 535 tokens to deliver 95.5\% recall, whereas ManiGit consumes only 237 tokens for 86.4\% recall, saving substantial API compute and avoiding the ``Lost-in-the-Middle'' degradation documented by Liu et al.~\cite{liu2024lost}.

Memory should be an unbounded, persistent, and versioned repository; Context should be a compact, highly curated, query-specific operative prompt compiled under a strict token knapsack.

\subsection{Contradiction Supersession vs. Ghost Context}
Our ablation experiments quantitatively isolate the threat of \textit{ghost context}. When an older fact (e.g., PostgreSQL) is contradicted by a newer fact (e.g., MongoDB), leaving both facts active in retrieval results in a 50\% failure rate on migration queries.

Probabilistic prompt instructions (e.g., ``Always prefer the most recent statement'') are insufficient when prompts contain dozens of competing chunks. Language models frequently blend conflicting statements or hallucinate compromises. ManiGit proves that contradiction resolution must occur \textbf{at the memory architecture level}: superseded facts must be explicitly archived so they are quarantined from the context compiler.

\subsection{Automated Infrastructure vs. Agent Tool Calling}
Comparing ManiGit to emerging Git-inspired agent tools (such as Git Context Controller~\cite{wu2025gcc} and GitOfThoughts~\cite{shekar2025gitofthoughts}) highlights an essential design choice: \textbf{should memory management be an agent-facing tool or an automated infrastructure substrate?}

When memory management is implemented as tool commands (\texttt{COMMIT}, \texttt{BRANCH}, \texttt{MERGE}), the agent must expend cognitive effort, prompt tokens, and multi-turn execution cycles deciding when and how to manage its memory. If the model forgets to commit a milestone, state is lost. ManiGit operates as a transparent system substrate: raw messages are automatically parsed, classified, committed, and compiled in the background. The agent interacts with users naturally, while the underlying repository guarantees immutability, provenance, and budget compliance.
